\section{基数}

\subsection{等势}

\begin{definition}[等势]
    设 $A, B$ 为集合，若存在 $A \to B$ 的双射，那么称 $A, B$ \textbf{等势}，记作 $A \sim B$。
\end{definition}

可以发现，等势关系是等价关系，具有自反性、对称性和传递性。

\begin{definition}[有穷集]
    设 $A$ 是集合，若 $\exists\,n \in \NN: A \sim n$，那么称 $A$ 为\textbf{有穷集}，否则称 $A$ 为\textbf{无穷集}。
\end{definition}

有了等势的概念，就可以推广集合“元素个数”的概念，定义集合的基数。

\begin{definition}[基数]
    设 $A$ 是集合，且 $n \in \NN,\ A \sim n$，那么称 $A$ 的\textbf{基数}为 $n$，记作 $\# A$ 或 $\operatorname{card}(A)$ 或 $|A|$。

    特别地，无穷集的基数使用特殊的记号，如 $\# \NN = \aleph_0$。
\end{definition}
